% ============================================================
%  يُترجَم بـ XeLaTeX  (في Overleaf: Menu → Compiler → XeLaTeX)
%  الخط: Amiri (متوفر في Overleaf، أو حمّله من github.com/aliftype/amiri)
% ============================================================
\documentclass[12pt,a4paper]{article}
\usepackage[a4paper,margin=1.7cm,top=2.3cm,bottom=2.9cm,headheight=15pt,footskip=1.5cm]{geometry}
\usepackage{amsmath,amssymb}
\usepackage{icomma}
\usepackage[english,bidi=default,layout=lists]{babel}
\babelprovide[import,main]{arabic}
\babelfont{rm}[Scale=1.08]{Amiri}
\babelfont{sf}[Scale=1.08]{Amiri}
\usepackage[most]{tcolorbox}
\usepackage{tikz}
\usetikzlibrary{arrows.meta,calc,shadings}
\usepackage{array}
\usepackage{colortbl}
\usepackage{fancyhdr}
\usepackage{titlesec}
\usepackage{enumitem}
\usepackage{collcell}
\usepackage{needspace}
\usepackage{etoolbox}
\newcolumntype{C}{>{\collectcell\ar}c<{\endcollectcell}}
\pretocmd{\section}{\Needspace{15\baselineskip}}{}{}
\pretocmd{\subsection}{\Needspace{11\baselineskip}}{}{}

\definecolor{cDef}{HTML}{1F4E9C}
\definecolor{cThm}{HTML}{B3261E}
\definecolor{cRem}{HTML}{2E7D32}
\definecolor{cEx}{HTML}{C76B00}
\definecolor{cMeth}{HTML}{6A1B9A}
\definecolor{cSol}{HTML}{00695C}
\definecolor{cSec}{HTML}{1F3A68}

\titleformat{\section}{\Large\bfseries\color{cSec}}{\thesection.}{0.6em}{}[{\color{cSec!60}\titlerule[1pt]}]
\titleformat{\subsection}{\large\bfseries\color{cThm!85!black}}{\thesubsection.}{0.5em}{}
\titlespacing*{\section}{0pt}{14pt}{8pt}
\titlespacing*{\subsection}{0pt}{10pt}{5pt}
\setlength{\parindent}{0pt}
\setlength{\parskip}{4pt}
\setlist{itemsep=2pt,topsep=3pt}

\newtcolorbox{arbox}[3][]{enhanced,
  colback=#2!5,colframe=#2!80!black,arc=2mm,boxrule=0.8pt,
  left=6pt,right=6pt,top=8pt,bottom=5pt,
  fonttitle=\bfseries,coltitle=white,
  title={\foreignlanguage{arabic}{#3}},
  attach boxed title to top right={xshift=-5mm,yshift=-3mm},
  boxed title style={colback=#2!80!black,arc=1.5mm,boxrule=0pt},
  before upper={\selectlanguage{arabic}},
  before skip=8pt,after skip=8pt,breakable,#1}
\newenvironment{BOX}[2]{\begin{otherlanguage}{english}\begin{arbox}{#1}{#2}}{\end{arbox}\end{otherlanguage}}
\newenvironment{tarif}[1][تعريف]{\begin{BOX}{cDef}{#1}}{\end{BOX}}
\newenvironment{khassia}[1][خاصية]{\begin{BOX}{cThm}{#1}}{\end{BOX}}
\newenvironment{molahada}[1][ملاحظة]{\begin{BOX}{cRem}{#1}}{\end{BOX}}
\newenvironment{mithal}[1][مثال]{\begin{BOX}{cEx}{#1}}{\end{BOX}}
\newenvironment{tariqa}[1][طريقة]{\begin{BOX}{cMeth}{#1}}{\end{BOX}}

\newcounter{exo}
\newenvironment{tamrin}[1]{\stepcounter{exo}\begin{BOX}{cSec}{التمرين \theexo: #1}}{\end{BOX}}
\newenvironment{hal}{\Needspace{12\baselineskip}\begin{otherlanguage}{english}\begin{arbox}[breakable]{cSol}{حل التمرين \theexo}}{\end{arbox}\end{otherlanguage}}

\newenvironment{LTR}{\begin{otherlanguage}{english}}{\end{otherlanguage}}
\newcommand{\ar}[1]{\foreignlanguage{arabic}{#1}}

\renewcommand{\dots}{\mathinner{\ldotp\ldotp\ldotp}}
\renewcommand{\ldots}{\mathinner{\ldotp\ldotp\ldotp}}
\renewcommand{\cdots}{\mathinner{\cdotp\cdotp\cdotp}}
\newcommand{\N}{\mathbb{N}}
\newcommand{\Z}{\mathbb{Z}}
\newcommand{\D}{\mathbb{D}}
\newcommand{\Q}{\mathbb{Q}}
\newcommand{\R}{\mathbb{R}}
\newcommand{\pgcd}{\operatorname{PGCD}}
\newcommand{\ppcm}{\operatorname{PPCM}}

% \NLseg{a}{b}{[}{]} ; \NLright{a}{[} ; \NLleft{b}{]}  (interval on a number line)
\newcommand{\NLseg}[4]{\begin{LTR}\begin{tikzpicture}[baseline=-0.6ex]
 \draw[-{Stealth[length=2mm]},gray!70](0,0)--(3.4,0);
 \draw[line width=2pt,cThm!80!black](0.7,0)--(2.7,0);
 \node[cThm!80!black,font=\LARGE] at (0.7,0){$#3$};
 \node[cThm!80!black,font=\LARGE] at (2.7,0){$#4$};
 \node[below=3pt,font=\footnotesize] at (0.7,0){$#1$};
 \node[below=3pt,font=\footnotesize] at (2.7,0){$#2$};
\end{tikzpicture}\end{LTR}}
\newcommand{\NLright}[2]{\begin{LTR}\begin{tikzpicture}[baseline=-0.6ex]
 \draw[-{Stealth[length=2mm]},gray!70](0,0)--(3.4,0);
 \draw[line width=2pt,cThm!80!black](1.3,0)--(3.3,0);
 \node[cThm!80!black,font=\LARGE] at (1.3,0){$#2$};
 \node[below=3pt,font=\footnotesize] at (1.3,0){$#1$};
\end{tikzpicture}\end{LTR}}
\newcommand{\NLleft}[2]{\begin{LTR}\begin{tikzpicture}[baseline=-0.6ex]
 \draw[-{Stealth[length=2mm]},gray!70](0,0)--(3.4,0);
 \draw[line width=2pt,cThm!80!black](0.05,0)--(2.1,0);
 \node[cThm!80!black,font=\LARGE] at (2.1,0){$#2$};
 \node[below=3pt,font=\footnotesize] at (2.1,0){$#1$};
\end{tikzpicture}\end{LTR}}

% ---------- social-media icons (pure TikZ) ----------
\definecolor{cFB}{HTML}{1877F2}
\definecolor{cTG}{HTML}{229ED9}
\definecolor{cIG1}{HTML}{FEDA75}
\definecolor{cIG2}{HTML}{D62976}
\definecolor{cIG3}{HTML}{4F5BD5}
\newcommand{\fbicon}{\begin{LTR}\begin{tikzpicture}[baseline=-0.75ex,x=1em,y=1em]
 \fill[cFB] (0,0) circle (0.55);
 \node[white,inner sep=0,font=\usefont{TU}{lmss}{bx}{n}\small] at (0.04,-0.02) {f};
\end{tikzpicture}\end{LTR}}
\newcommand{\igicon}{\begin{LTR}\begin{tikzpicture}[baseline=-0.75ex,x=1em,y=1em]
 \shade[lower left=cIG1,lower right=cIG2,upper left=cIG2,upper right=cIG3,rounded corners=0.28em] (-0.55,-0.55) rectangle (0.55,0.55);
 \draw[white,line width=0.09em,rounded corners=0.18em] (-0.36,-0.36) rectangle (0.36,0.36);
 \draw[white,line width=0.09em] (0,0) circle (0.17);
 \fill[white] (0.22,0.22) circle (0.05);
\end{tikzpicture}\end{LTR}}
\newcommand{\tgicon}{\begin{LTR}\begin{tikzpicture}[baseline=-0.75ex,x=1em,y=1em]
 \fill[cTG] (0,0) circle (0.55);
 \fill[white] (-0.33,0.0) -- (0.32,0.27) -- (0.2,-0.3) -- (0.0,-0.15) -- (-0.08,-0.28) -- (-0.08,-0.08) -- cycle;
 \fill[cTG!70!black] (-0.08,-0.08) -- (0.25,0.2) -- (0.0,-0.15) -- (-0.08,-0.28) -- cycle;
\end{tikzpicture}\end{LTR}}
% page number: black box, white two-digit number, at the right of the social row
\newcommand{\PageBox}{\begin{LTR}\tikz[baseline=(p.base)]\node[fill=black,text=white,inner sep=2.5pt,minimum width=1.6em,minimum height=1.3em,font=\usefont{TU}{lmss}{bx}{n}\small](p){\ifnum\value{page}<10 0\fi\arabic{page}};\end{LTR}}
% visual left-to-right order: Telegram | Instagram | Facebook | page box
\newcommand{\SocialRow}{\begin{otherlanguage}{english}%
 \ar{\textbf{الأستاذ سوايسية هشام للرياضيات}}\ \tgicon\hfill
 \textbf{prof\_hicham\_math}\ \igicon\hfill
 \ar{\textbf{الأستاذ سوايسية هشام للرياضيات}}\ \fbicon\hspace{2.2em}\PageBox\end{otherlanguage}}

\pagestyle{fancy}
\fancyhf{}
\renewcommand{\headrulewidth}{0.6pt}
\renewcommand{\footrulewidth}{0.4pt}
% \SetHeader{right head: document title}{left head: level -- stream}
\newcommand{\SetHeader}[2]{%
 \fancyhead[C]{\parbox{\textwidth}{\ar{#1\hfill #2}}}%
 \fancyfoot[C]{\PageBox}}
\usepackage{tkz-tab}
\usepackage{multicol}
% variation / sign table (LTR)
\newenvironment{TAB}{\par\begin{center}\begin{LTR}\begin{tikzpicture}[scale=0.95,every node/.style={font=\small}]}{\end{tikzpicture}\end{LTR}\end{center}}
% grid + axes, clipped: \begin{GR}{xmin}{xmax}{ymin}{ymax}{unit cm}
\newenvironment{GR}[5]{\par\begin{center}\begin{LTR}\begin{tikzpicture}[x=#5cm,y=#5cm,font=\footnotesize]
 \draw[gray!25,very thin] (#1,#3) grid (#2,#4);
 \draw[-{Stealth[length=2mm]}] (#1,0)--(#2+0.3,0) node[below left]{$x$};
 \draw[-{Stealth[length=2mm]}] (0,#3)--(0,#4+0.3) node[below left]{$y$};
 \node[below left] at (0,0){$O$};
 \begin{scope}\clip (#1,#3) rectangle (#2,#4);}
{\end{scope}\end{tikzpicture}\end{LTR}\end{center}}
\newcommand{\Ticks}[2]{\foreach \t in {#1} \node[below,font=\scriptsize] at (\t,0){$\t$}; \foreach \t in {#2} \node[left,font=\scriptsize] at (0,\t){$\t$};}
\newcommand{\bx}[1]{\ensuremath{\colorbox{cDef!12}{$\displaystyle #1$}}}
\newcommand{\TwoCol}[4]{\par\begin{LTR}\noindent\begin{minipage}[t]{#1\textwidth}\selectlanguage{arabic}#2\end{minipage}\hfill\begin{minipage}[t]{#3\textwidth}\selectlanguage{arabic}#4\end{minipage}\end{LTR}\par}
\newcommand{\RL}[4]{\TwoCol{#3}{#4}{#1}{#2}}  % \RL{w1}{right content}{w2}{left content}
\newcommand{\ok}{$\textcolor{cRem}{\checkmark}$}
\newcommand{\limn}{\displaystyle\lim_{n\to+\infty}}
% sub-questions, each on its own line
\newenvironment{sq}{\begin{itemize}[leftmargin=2.4em,itemsep=3pt]}{\end{itemize}}
\newcommand{\qa}{\item[(أ)]}\newcommand{\qb}{\item[(ب)]}\newcommand{\qc}{\item[(جـ)]}\newcommand{\qd}{\item[(د)]}\newcommand{\qe}{\item[(هـ)]}
\newcommand{\pa}{\par\smallskip\textbf{(أ)} }\newcommand{\pb}{\par\smallskip\textbf{(ب)} }\newcommand{\pc}{\par\smallskip\textbf{(جـ)} }\newcommand{\pd}{\par\smallskip\textbf{(د)} }\newcommand{\pe}{\par\smallskip\textbf{(هـ)} }
\newcommand{\QQ}[1]{\par\smallskip\textbf{#1.} }
\newenvironment{tmr}{\stepcounter{exo}\begin{BOX}{cSec}{تمرين محلول \theexo}}{\end{BOX}}
\newcommand{\ABC}[3]{\\[2pt]\begin{tabular}{p{0.28\textwidth}p{0.28\textwidth}p{0.28\textwidth}}(جـ) \ #3 & (ب) \ #2 & (أ) \ #1\end{tabular}}

\newcommand{\nt}[1]{\hfill$\textcolor{cThm}{\mathbf{(#1)}}$}
\SetHeader{\small\bfseries\color{cSec}فرض المراقبة الذاتية رقم 01}{\small\color{cSec}السنة الثانية ثانوي -- تقني رياضي ورياضيات}
\begin{document}
\begin{center}
{\large\bfseries الجمهورية الجزائرية الديمقراطية الشعبية}\\[2pt]
{\bfseries وزارة التربية الوطنية}\\[2pt]
{\bfseries الديوان الوطني للتعليم والتكوين عن بعد}\\[2pt]
{\bfseries المركز الولائي سوق أهراس}
\end{center}
\noindent\textbf{السنة الدراسية: 2026/2027}\hfill\textbf{المستوى: سنة ثانية ثانوي شعبة تقني رياضي ورياضيات}
\par\smallskip
\begin{LTR}
\begin{tcolorbox}[enhanced,colback=cSec,colframe=cSec,arc=2mm,halign=center,boxrule=0pt,top=6pt,bottom=6pt,before upper={\selectlanguage{arabic}}]
{\color{white}\Large\bfseries فرض المراقبة الذاتية رقم 01 -- الموضوع الثالث}
\end{tcolorbox}
\end{LTR}

\begin{tamrin}{كثيرات الحدود (06 نقاط)}
نعتبر كثير الحدود $P$ المعرّف على $\R$ بـ: $P(x)=x^3-2x^2-5x+6$.
\QQ{1} احسب $P(1)$ ، ماذا تستنتج؟
\QQ{2} باستعمال جدول هورنر ، عيّن كثير الحدود $Q(x)$ بحيث: $P(x)=(x-1)Q(x)$.
\QQ{3} حلّل $P(x)$ إلى جداء عوامل من الدرجة الأولى ، ثم حُلّ في $\R$ المعادلة: $P(x)=0$.
\QQ{4} ادرس إشارة $P(x)$ ، ثم استنتج حلول المتراجحة: $P(x)>0$.
\QQ{5} استنتج في $\R$ حلول المعادلة: $(x+1)^3-2(x+1)^2-5(x+1)+6=0$.
\end{tamrin}

\begin{tamrin}{الاشتقاقية (07 نقاط)}
$f$ الدالة المعرّفة على $\R$ بـ: $f(x)=x^3-3x^2+4$ ، و $(C_f)$ تمثيلها البياني في معلم متعامد ومتجانس $\left(O;\vec i,\vec j\right)$.
\QQ{1} باستعمال تعريف العدد المشتق ، بيّن أن $f$ قابلة للاشتقاق عند $1$ وأن: $f^{\prime}(1)=-3$.
\QQ{2} اكتب معادلة للمماس $(T)$ للمنحنى $(C_f)$ عند النقطة ذات الفاصلة $1$ ، ثم استنتج قيمة تقريبية للعدد $f(1{,}002)$.
\QQ{3} احسب $f^{\prime}(x)$ ، ثم ادرس اتجاه تغيّر الدالة $f$ وشكّل جدول تغيّراتها (دون حساب النهايات) ، وعيّن قيمها الحدية المحلية.
\QQ{4} تحقّق أنه من أجل كل عدد حقيقي $x$: $f(x)=(x-2)^2(x+1)$ ، ثم ادرس وضعية $(C_f)$ بالنسبة إلى محور الفواصل.
\QQ{5} بيّن أنه من أجل كل عدد حقيقي $x$: $f(x)-(-3x+5)=(x-1)^3$ ، ثم ادرس وضعية $(C_f)$ بالنسبة إلى المماس $(T)$.
\QQ{6} عيّن فواصل النقط من $(C_f)$ التي يكون فيها المماس موازياً للمستقيم ذي المعادلة $y=9x$.
\end{tamrin}

\begin{tamrin}{الاحتمالات (07 نقاط)}
يحتوي صندوق على $5$ قريصات متماثلة لا نفرّق بينها باللمس: ثلاث حمراء تحمل الأرقام $1$ ، $2$ ، $3$ ، واثنتان خضراوان تحملان الرقمين $1$ ، $2$. نسحب عشوائياً \textbf{في آن واحد} قريصتين من الصندوق.
\QQ{1} مثّل مجموعة الإمكانيات بجدول ذي مدخلين ، ثم عيّن عددها.
\QQ{2} احسب احتمال كل حادثة من الحوادث التالية:
\begin{sq}
\qa $A$: «القريصتان من نفس اللون».
\qb $B$: «مجموع الرقمين يساوي $3$».
\qc $C$: «سحب قريصة خضراء على الأقل».
\qd $A\cup B$.
\end{sq}
\QQ{3} ليكن $X$ المتغير العشوائي الذي يرفق بكل سحبة مجموع الرقمين المسجلين على القريصتين.
\begin{sq}
\qa عيّن القيم الممكنة للمتغير العشوائي $X$ ، ثم عرّف قانون احتماله.
\qb احسب الأمل الرياضياتي $E(X)$ ، والتباين $V(X)$ ، والانحراف المعياري $\sigma(X)$.
\end{sq}
\end{tamrin}
\begin{center}\textbf{\color{cSec}بالتوفيق}\end{center}

\newpage
\begin{LTR}
\begin{tcolorbox}[enhanced,colback=cSol,colframe=cSol,arc=2mm,halign=center,boxrule=0pt,top=6pt,bottom=6pt,before upper={\selectlanguage{arabic}}]
{\color{white}\Large\bfseries الحل النموذجي وسلم التنقيط}
\end{tcolorbox}
\end{LTR}
\Needspace{12\baselineskip}\begin{BOX}{cSol}{حل التمرين 1}
\QQ{1} $P(1)=1-2-5+6=\bx{0}$ ، إذن $1$ جذر لـ $P(x)$ ، و $P(x)$ يقبل القسمة على $(x-1)$.\hfill$\textcolor{cThm}{\mathbf{(0,5)}}$
\QQ{2} جدول هورنر (المعاملات $1$ ، $-2$ ، $-5$ ، $6$ والجذر $1$):
\begin{center}\begin{LTR}{\renewcommand{\arraystretch}{1.5}\begin{tabular}{|c|c|c|c|c|}\hline
\rowcolor{cSec!12} & $1$ & $-2$ & $-5$ & $6$\\\hline
$1$ & & $1$ & $-1$ & $-6$\\\hline
 & $1$ & $-1$ & $-6$ & $0$\\\hline
\end{tabular}}\end{LTR}\end{center}
إذن: \bx{Q(x)=x^2-x-6} و $P(x)=(x-1)\left(x^2-x-6\right)$.\hfill$\textcolor{cThm}{\mathbf{(1,5)}}$
\QQ{3} $x^2-x-6$: $\Delta=1+24=25$ ، $x_1=\dfrac{1-5}2=-2$ ، $x_2=\dfrac{1+5}2=3$ ، ومنه $x^2-x-6=(x+2)(x-3)$.
\par \bx{P(x)=(x-1)(x+2)(x-3)} ، ومجموعة حلول $P(x)=0$: \bx{S=\{-2;1;3\}}\hfill$\textcolor{cThm}{\mathbf{(1,5)}}$
\QQ{4} جدول الإشارة:
\begin{TAB}
\tkzTabInit[lgt=2.2,espcl=1.6]{$x$/0.8,$x+2$/0.7,$x-1$/0.7,$x-3$/0.7,$P(x)$/0.8}{$-\infty$,$-2$,$1$,$3$,$+\infty$}
\tkzTabLine{,-,z,+,t,+,t,+,}
\tkzTabLine{,-,t,-,z,+,t,+,}
\tkzTabLine{,-,t,-,t,-,z,+,}
\tkzTabLine{,-,z,+,z,-,z,+,}
\end{TAB}
حلول المتراجحة $P(x)>0$: \bx{S=]-2;1[\cup]3;+\infty[}\hfill$\textcolor{cThm}{\mathbf{(1,5)}}$
\QQ{5} المعادلة تكتب $P(x+1)=0$ ، أي $x+1\in\{-2;1;3\}$ ، ومنه $x\in\{-3;0;2\}$: \bx{S=\{-3;0;2\}}\hfill$\textcolor{cThm}{\mathbf{(1)}}$
\end{BOX}

\Needspace{12\baselineskip}\begin{BOX}{cSol}{حل التمرين 2}
\QQ{1} $f(1)=2$. من أجل $h\neq0$: $f(1+h)=1+3h+3h^2+h^3-3-6h-3h^2+4=2-3h+h^3$.
\par $\dfrac{f(1+h)-f(1)}h=\dfrac{-3h+h^3}h=-3+h^2$ ، و $\lim\limits_{h\to0}\left(-3+h^2\right)=-3$ ، إذن $f$ قابلة للاشتقاق عند $1$ و \bx{f^{\prime}(1)=-3}\hfill$\textcolor{cThm}{\mathbf{(1,5)}}$
\QQ{2} $(T):y=f^{\prime}(1)(x-1)+f(1)=-3(x-1)+2$ ، أي \bx{(T):y=-3x+5}
\par التقريب التآلفي: $f(1{,}002)\approx-3(1{,}002)+5=\bx{1{,}994}$\hfill$\textcolor{cThm}{\mathbf{(1)}}$
\QQ{3} \bx{f^{\prime}(x)=3x^2-6x=3x(x-2)} ؛ $f(0)=4$ ، $f(2)=8-12+4=0$.
\begin{TAB}
\tkzTabInit[lgt=1.6,espcl=2.2]{$x$/0.8,$f^{\prime}(x)$/0.8,$f(x)$/1.6}{$-\infty$,$0$,$2$,$+\infty$}
\tkzTabLine{,+,z,-,z,+,}
\tkzTabVar{-/ ,+/$4$,-/$0$,+/ }
\end{TAB}
$f$ متزايدة تماماً على $]-\infty;0]$ و $[2;+\infty[$ ، ومتناقصة تماماً على $[0;2]$. قيمة عظمى محلية \bx{f(0)=4} ، وقيمة صغرى محلية \bx{f(2)=0}\hfill$\textcolor{cThm}{\mathbf{(1,5)}}$
\QQ{4} $(x-2)^2(x+1)=\left(x^2-4x+4\right)(x+1)=x^3-3x^2+4=f(x)$.
\par $(x-2)^2\geqslant0$ ، فإشارة $f(x)$ من إشارة $x+1$:
\begin{sq}
\item $x<-1$: $f(x)<0$ و $(C_f)$ \textbf{تحت} محور الفواصل.
\item $x>-1$ و $x\neq2$: $f(x)>0$ و $(C_f)$ \textbf{فوق} محور الفواصل.
\item $(C_f)$ يقطع محور الفواصل في $(-1;0)$ ويمسّه في $(2;0)$.
\end{sq}\hfill$\textcolor{cThm}{\mathbf{(1)}}$
\QQ{5} $f(x)-(-3x+5)=x^3-3x^2+3x-1=\bx{(x-1)^3}$
\par $x<1$: $(x-1)^3<0$ و $(C_f)$ \textbf{تحت} $(T)$ ؛ $x>1$: $(C_f)$ \textbf{فوق} $(T)$ ؛ و $(C_f)$ يقطع $(T)$ في النقطة $(1;2)$.\hfill$\textcolor{cThm}{\mathbf{(1)}}$
\QQ{6} $f^{\prime}(x)=9$ معناه $3x^2-6x-9=0$ أي $x^2-2x-3=0$ ، أي $(x-3)(x+1)=0$: الفاصلتان \bx{x=-1} و \bx{x=3}\hfill$\textcolor{cThm}{\mathbf{(1)}}$
\end{BOX}

\Needspace{12\baselineskip}\begin{BOX}{cSol}{حل التمرين 3}
\QQ{1} نرمز للقريصات: $R_1$ ، $R_2$ ، $R_3$ ، $V_1$ ، $V_2$. السحب في آن واحد: الترتيب غير مهم ولا تكرار ، فنستعمل نصف الجدول فوق القطر. في كل خانة نكتب مجموع الرقمين (قيمة $X$):
\begin{center}\begin{LTR}{\renewcommand{\arraystretch}{1.5}\begin{tabular}{|c|c|c|c|c|c|}\hline
\rowcolor{cSec!12}$\text{I}\backslash\text{II}$ & $R_1$ & $R_2$ & $R_3$ & $V_1$ & $V_2$\\\hline
$R_1$ & \cellcolor{gray!25} & \cellcolor{cEx!35}$3$ & \cellcolor{cEx!35}$4$ & $2$ & $3$\\\hline
$R_2$ & \cellcolor{gray!25} & \cellcolor{gray!25} & \cellcolor{cEx!35}$5$ & $3$ & $4$\\\hline
$R_3$ & \cellcolor{gray!25} & \cellcolor{gray!25} & \cellcolor{gray!25} & $4$ & $5$\\\hline
$V_1$ & \cellcolor{gray!25} & \cellcolor{gray!25} & \cellcolor{gray!25} & \cellcolor{gray!25} & \cellcolor{cEx!35}$3$\\\hline
$V_2$ & \cellcolor{gray!25} & \cellcolor{gray!25} & \cellcolor{gray!25} & \cellcolor{gray!25} & \cellcolor{gray!25}\\\hline
\end{tabular}}\end{LTR}\end{center}
عدد الإمكانيات: \bx{\dfrac{5\times4}2=10} ، وهي متساوية الاحتمال. (البرتقالي: الحادثة $A$.)\hfill$\textcolor{cThm}{\mathbf{(1,5)}}$
\QQ{2}
\pa $A$: $\{R_1;R_2\}$ ، $\{R_1;R_3\}$ ، $\{R_2;R_3\}$ ، $\{V_1;V_2\}$: \bx{P(A)=\tfrac4{10}=\tfrac25}\hfill$\textcolor{cThm}{\mathbf{(0,5)}}$
\pb $B$: الخانات التي فيها $3$: $\{R_1;R_2\}$ ، $\{R_1;V_2\}$ ، $\{R_2;V_1\}$ ، $\{V_1;V_2\}$: \bx{P(B)=\tfrac4{10}=\tfrac25}\hfill$\textcolor{cThm}{\mathbf{(0,5)}}$
\pc $\overline C$: «القريصتان حمراوان» ، $P\left(\overline C\right)=\frac3{10}$ ، إذن \bx{P(C)=1-\tfrac3{10}=\tfrac7{10}}\hfill$\textcolor{cThm}{\mathbf{(0,5)}}$
\pd $A\cap B=\left\{\{R_1;R_2\};\{V_1;V_2\}\right\}$ ، $P(A\cap B)=\frac2{10}$: $P(A\cup B)=\dfrac4{10}+\dfrac4{10}-\dfrac2{10}=\bx{\tfrac6{10}=\tfrac35}$\hfill$\textcolor{cThm}{\mathbf{(1)}}$
\QQ{3}
\pa $X\in\{2;3;4;5\}$. بعدّ الخانات:
\begin{center}\begin{LTR}{\renewcommand{\arraystretch}{1.6}\begin{tabular}{|C|c|c|c|c|C|}\hline
\rowcolor{cSec!12}$x_i$ & $2$ & $3$ & $4$ & $5$ & المجموع\\\hline
$P(X=x_i)$ & $\frac1{10}$ & $\frac4{10}$ & $\frac3{10}$ & $\frac2{10}$ & $1$\\\hline
\end{tabular}}\end{LTR}\end{center}
\hfill$\textcolor{cThm}{\mathbf{(1,5)}}$
\pb $E(X)=\dfrac{2+12+12+10}{10}=\bx{\tfrac{18}5=3{,}6}$\hfill$\textcolor{cThm}{\mathbf{(0,5)}}$
\par $V(X)=\dfrac{4+36+48+50}{10}-(3{,}6)^2=13{,}8-12{,}96=\bx{0{,}84}$ ؛ $\sigma(X)=\sqrt{0{,}84}\approx\bx{0{,}92}$\hfill$\textcolor{cThm}{\mathbf{(1)}}$
\end{BOX}

\begin{BOX}{cSec}{سلم التنقيط}
\begin{center}\begin{LTR}{\renewcommand{\arraystretch}{1.5}\begin{tabular}{|C|c|c|c|c|c|c|C|}\hline
\rowcolor{cSec!12}التمرين & \ar{س1} & \ar{س2} & \ar{س3} & \ar{س4} & \ar{س5} & \ar{س6} & المجموع\\\hline
الأول & $0{,}5$ & $1{,}5$ & $1{,}5$ & $1{,}5$ & $1$ & -- & $06$\\\hline
الثاني & $1{,}5$ & $1$ & $1{,}5$ & $1$ & $1$ & $1$ & $07$\\\hline
الثالث & $1{,}5$ & $2{,}5$ & $3$ & -- & -- & -- & $07$\\\hline
\rowcolor{cSec!12}المجموع & \multicolumn{6}{c|}{} & $20$\\\hline
\end{tabular}}\end{LTR}\end{center}
\end{BOX}
\end{document}
